\section{Additional Simulation Results}
\label{app:results}

\subsection{Per-Task Results}
\label{app:per_task}
% Wrap each Delta cell as \dcell{value}{text} once results arrive (shading rule in the caption).
% Transcribed from results/realtime_c/ ACT/DP per-task CSV and reports (INFER0); other methods pending.
\begin{table*}[t]
\centering\footnotesize
\caption{\tablelang{\textbf{Per-task held-out success} (\%) on \bench for every method of Table~\mainref{tab:main}, mean of
three training seeds for our variants and one for baselines, real-time protocol. DP: Diffusion Policy. DVLA: DynamicVLA. DrVLA: DreamVLA. OWAM: OpenWAM-$\alpha$.
DWAM: DynamicWAM. +: with our SE(3) history tokens. $\sigma$: \ours with
nominal synthetic correlated pose noise (Appendix~\ref{app:pose_noise}). $\Delta$: \ours minus the strongest other method on that task, with its 95\% paired interval. The last two rows give the position error at first contact $e_p$ (mm), averaged over the 13 tasks with a contact, and the number $n$ of episodes with a contact (Appendix~\ref{app:metrics}); --: no 13-task mean. Task codes are
listed in Table~\ref{tab:roles}. Shading of $\Delta$: blue for higher SR, orange for lower, saturating at 15 points.}{\textbf{\bench 各任务的留出成功率}（\%），涵盖表~\mainref{tab:main} 中的全部方法，本文各变体取三个训练种子的平均、基线训练一次，实时协议。
DP：Diffusion Policy。DVLA：DynamicVLA。DrVLA：DreamVLA。OWAM：OpenWAM-$\alpha$。DWAM：DynamicWAM。+：加入本文 SE(3) 历史 token。
$\sigma$：使用标称合成相关位姿噪声的 \ours（附录~\ref{app:pose_noise}）。$\Delta$：该任务上 \ours 相对最强其他方法之差及其 95\% 配对区间。最后两行为首次接触时的位置误差 $e_p$（mm，在有接触的 13 个任务上平均）与发生接触的回合数 $n$（附录~\ref{app:metrics}）；--：无 13 任务均值。
任务代号见表~\ref{tab:roles}。$\Delta$ 的底色：蓝色表示 SR 更高，橙色表示更低，在 15 个百分点处饱和。}}
\label{tab:supp_per_task}
\setlength{\tabcolsep}{2.4pt}
\begin{tabular}{@{}lccccccccccccc@{}}
\toprule
\tablelang{Task}{任务} & ACT & DP & $\pi_{0.5}$ & $\pi_{0.5}$+ & DVLA & PUMA & DrVLA & OWAM & DWAM & DWAM+ & $\sigma$ & \textbf{\tablelang{Ours}{本文}} & $\Delta$ [95\% CI] \\
\midrule
\tablelang{Harmonic grasp}{谐波抓取} & 48.0 & 0.0 & \resultplaceholder & \resultplaceholder & \resultplaceholder & \resultplaceholder & \resultplaceholder & \resultplaceholder & \resultplaceholder & \resultplaceholder & \resultplaceholder & \resultplaceholder & \resultplaceholder[{+00.0 [00.0, 00.0]}] \\
\tablelang{Conveyor grasp}{传送带抓取} & 46.0 & 0.0 & \resultplaceholder & \resultplaceholder & \resultplaceholder & \resultplaceholder & \resultplaceholder & \resultplaceholder & \resultplaceholder & \resultplaceholder & \resultplaceholder & \resultplaceholder & \resultplaceholder[{+00.0 [00.0, 00.0]}] \\
\tablelang{Translating grasp}{平移抓取} & 18.0 & 0.0 & \resultplaceholder & \resultplaceholder & \resultplaceholder & \resultplaceholder & \resultplaceholder & \resultplaceholder & \resultplaceholder & \resultplaceholder & \resultplaceholder & \resultplaceholder & \resultplaceholder[{+00.0 [00.0, 00.0]}] \\
\tablelang{Rotating grasp}{旋转抓取} & 21.0 & 0.0 & \resultplaceholder & \resultplaceholder & \resultplaceholder & \resultplaceholder & \resultplaceholder & \resultplaceholder & \resultplaceholder & \resultplaceholder & \resultplaceholder & \resultplaceholder & \resultplaceholder[{+00.0 [00.0, 00.0]}] \\
\tablelang{Occluded grasp}{遮挡抓取} & 32.0 & 0.0 & \resultplaceholder & \resultplaceholder & \resultplaceholder & \resultplaceholder & \resultplaceholder & \resultplaceholder & \resultplaceholder & \resultplaceholder & \resultplaceholder & \resultplaceholder & \resultplaceholder[{+00.0 [00.0, 00.0]}] \\
\midrule
\tablelang{Moving receptacle}{移动容器} & 26.0 & 0.0 & \resultplaceholder & \resultplaceholder & \resultplaceholder & \resultplaceholder & \resultplaceholder & \resultplaceholder & \resultplaceholder & \resultplaceholder & \resultplaceholder & \resultplaceholder & \resultplaceholder[{+00.0 [00.0, 00.0]}] \\
\tablelang{Rotating receptacle}{旋转容器} & 0.0 & 0.0 & \resultplaceholder & \resultplaceholder & \resultplaceholder & \resultplaceholder & \resultplaceholder & \resultplaceholder & \resultplaceholder & \resultplaceholder & \resultplaceholder & \resultplaceholder & \resultplaceholder[{+00.0 [00.0, 00.0]}] \\
\midrule
\tablelang{Wire passage}{导线穿越} & 99.0 & 0.0 & \resultplaceholder & \resultplaceholder & \resultplaceholder & \resultplaceholder & \resultplaceholder & \resultplaceholder & \resultplaceholder & \resultplaceholder & \resultplaceholder & \resultplaceholder & \resultplaceholder[{+00.0 [00.0, 00.0]}] \\
\tablelang{Wire passage, full}{导线穿越（完整）} & 0.0 & 0.0 & \resultplaceholder & \resultplaceholder & \resultplaceholder & \resultplaceholder & \resultplaceholder & \resultplaceholder & \resultplaceholder & \resultplaceholder & \resultplaceholder & \resultplaceholder & \resultplaceholder[{+00.0 [00.0, 00.0]}] \\
\midrule
\tablelang{Ball capture}{小球捕获} & 36.0 & 0.0 & \resultplaceholder & \resultplaceholder & \resultplaceholder & \resultplaceholder & \resultplaceholder & \resultplaceholder & \resultplaceholder & \resultplaceholder & \resultplaceholder & \resultplaceholder & \resultplaceholder[{+00.0 [00.0, 00.0]}] \\
\tablelang{Ball capture, table}{小球捕获（桌面）} & 71.0 & 0.0 & \resultplaceholder & \resultplaceholder & \resultplaceholder & \resultplaceholder & \resultplaceholder & \resultplaceholder & \resultplaceholder & \resultplaceholder & \resultplaceholder & \resultplaceholder & \resultplaceholder[{+00.0 [00.0, 00.0]}] \\
\tablelang{Slider capture}{滑块捕获} & 57.0 & 0.0 & \resultplaceholder & \resultplaceholder & \resultplaceholder & \resultplaceholder & \resultplaceholder & \resultplaceholder & \resultplaceholder & \resultplaceholder & \resultplaceholder & \resultplaceholder & \resultplaceholder[{+00.0 [00.0, 00.0]}] \\
\tablelang{Can capture}{罐体捕获} & 90.0 & 0.0 & \resultplaceholder & \resultplaceholder & \resultplaceholder & \resultplaceholder & \resultplaceholder & \resultplaceholder & \resultplaceholder & \resultplaceholder & \resultplaceholder & \resultplaceholder & \resultplaceholder[{+00.0 [00.0, 00.0]}] \\
\midrule
\tablelang{Post-impact grasp}{撞击后抓取} & 44.0 & 0.0 & \resultplaceholder & \resultplaceholder & \resultplaceholder & \resultplaceholder & \resultplaceholder & \resultplaceholder & \resultplaceholder & \resultplaceholder & \resultplaceholder & \resultplaceholder & \resultplaceholder[{+00.0 [00.0, 00.0]}] \\
\midrule
\tablelang{Task-balanced}{任务等权} & 42.0 & 0.0 & \resultplaceholder & \resultplaceholder & \resultplaceholder & \resultplaceholder & \resultplaceholder & \resultplaceholder & \resultplaceholder & \resultplaceholder & \resultplaceholder & \resultplaceholder & \resultplaceholder[{+00.0 [00.0, 00.0]}] \\
\midrule
$e_p$ (mm) & 32.2 & -- & \resultplaceholder & \resultplaceholder & \resultplaceholder & \resultplaceholder & \resultplaceholder & \resultplaceholder & \resultplaceholder & \resultplaceholder & \resultplaceholder & \resultplaceholder & -- \\
$n$ & 1109 & 6 & \resultplaceholder & \resultplaceholder & \resultplaceholder & \resultplaceholder & \resultplaceholder & \resultplaceholder & \resultplaceholder & \resultplaceholder & \resultplaceholder & \resultplaceholder & -- \\
\bottomrule
\end{tabular}
\end{table*}

Table~\ref{tab:supp_per_task} uses the Native real-time protocol with +0\,ms extra delay.
\ours reaches 93\% on moving-receptacle placement and 66\% on rotating-receptacle
placement, compared with 63\% and 11\% for $\pi_{0.5}$+SE(3).
On full wire passage, \ours reaches 40\% versus 12\% for DynamicVLA.
It ties the strongest
baseline on phase-only wire passage, tabletop ball capture and can capture.
The largest recorded failure classes for \ours are rim impulse violations
(160 of 1{,}400 scenes) and wire contacts (55 scenes).

\subsection{Latency}
\label{app:protocols}
Table~\ref{tab:supp_latency} lists all eleven configurations at five added
latencies; Fig.~\mainref{fig:latency} plots all 55 accepted points.


% Transcribed from results/realtime_c/latency-five-point-20261104.json: eleven configurations, 55 accepted points; one training seed each.
\begin{table}[H]
\centering\footnotesize
\caption{\tablelang{\textbf{Success versus added latency} on 1{,}400 held-out \bench scenes per point (SR in \%, real-time protocol). The 0 and
150\,ms columns equal the Native and +150\,ms columns of Table~\mainref{tab:main}. Blue/bold and light blue mark the best and second-best values per column.}{\textbf{成功率随附加延迟的变化}：每点使用 1{,}400 个 \bench 留出场景（SR，\%；实时协议）。0 与 150\,ms 两列分别等于表~\mainref{tab:main}
的原生与 +150\,ms 列。 粗体加蓝底与浅蓝底分别标记各列最佳与次佳。}}
\label{tab:supp_latency}
\setlength{\tabcolsep}{3pt}
\begin{tabular}{@{}lccccc@{}}
\toprule
 & \multicolumn{5}{c}{\tablelang{Added latency (ms)}{附加延迟（ms）}} \\
\cmidrule(lr){2-6}
\tablelang{Method}{方法} & 0 & 50 & 100 & 150 & 200 \\
\midrule
ACT & 41.6 & 43.5 & \secondcell{36.5} & 24.3 & 16.5 \\
Diffusion Policy & 30.7 & 26.1 & 21.0 & 15.9 & 14.3 \\
$\pi_{0.5}$ & 53.8 & \secondcell{47.2} & \bestcell{39.6} & 27.8 & 13.6 \\
$\pi_{0.5}$ + SE(3) & 52.8 & 39.8 & 32.4 & \secondcell{28.2} & 15.1 \\
DynamicVLA & 50.9 & 46.4 & 30.6 & 23.4 & 17.8 \\
PUMA & \secondcell{59.6} & \bestcell{53.2} & 36.4 & \bestcell{31.1} & \bestcell{27.1} \\
OpenWAM-$\alpha$ & 46.1 & 36.1 & 27.9 & 24.7 & \secondcell{22.1} \\
OpenWAM-$\alpha$ (off) & 15.8 & 13.7 & 14.2 & 10.4 & 9.5 \\
DynamicWAM & 40.5 & 27.6 & 25.0 & 20.6 & 20.9 \\
DynamicWAM + SE(3) & 36.8 & 30.6 & 23.9 & 21.9 & 21.3 \\
\midrule
\ours & \bestcell{64.3} & 43.4 & 25.5 & 9.9 & 6.6 \\
\bottomrule
\end{tabular}
\end{table}


\noindent\begin{minipage}{\linewidth}
\paragraph{Timing.}
\label{app:latency_results}
Table~\ref{tab:supp_timing} reports replay timings by stage for the development
checkpoint at 20{,}696 updates.

\begin{table}[H]
\centering\footnotesize
% Transcribed from results/realtime_c/stage-timing-20261104.json: accepted compiled BF16 C stage measurement, separate eager readouts and unprobed totals. Body only.
\caption{\tablelang{\textbf{Latency breakdown} of compiled BF16 \ours on one \cfgGPU{} at batch size 1 (ms, p50 / p95). Stage rows are CUDA-event spans over 900 queries; totals are uninstrumented wall time, excluding pose tracking. The original Table~\mainref{tab:main} calibration is retained alongside this run's control. Readouts are 900 separate BF16 eager forwards on one fixed recorded World feature; they are absent from deployment. Stage percentiles are not additive. Magenta dashes mark pending robot measurements.}{\textbf{延迟分解}（BF16编译版 \ours，\cfgGPU{}，batch 为1，ms；p50 / p95）。阶段项为900次查询的CUDA事件区间，合计为无探针墙钟时间，不含位姿跟踪；同时保留表~\mainref{tab:main} 的原始标定与本轮对照。仅训练读出在同一固定录制World特征上单独执行900次BF16 eager前向，不参与部署。阶段分位数不相加；品红色破折号表示待测的真机项。}}
\label{tab:supp_timing}
\setlength{\tabcolsep}{4pt}
\begin{tabular}{@{}lcc@{}}
\toprule
\tablelang{Stage}{阶段} & p50 & p95 \\
\midrule
\tablelang{Pose tracking (per frame, robot)}{位姿跟踪（每帧，真机）} & \resultplaceholder & \resultplaceholder \\
\tablelang{Observation encoding + H2D}{观测编码与传输} & 4.83 & 5.28 \\
\tablelang{World/KV preparation (once/query)}{World/KV准备（每次查询一次）} & 3.93 & 4.22 \\
\tablelang{Action generation (\cfgEuler{} Euler steps)}{动作生成（\cfgEuler{} 个Euler步）} & 4.27 & 4.41 \\
\tablelang{Policy total (Table 1 calibration)}{策略合计（表1标定）} & 14.59 & 15.11 \\
\tablelang{Policy total (unprobed control)}{策略合计（本轮无探针对照）} & 14.66 & 15.65 \\
\tablelang{Training-only readouts (separate, eager)}{仅训练读出（单独、eager）} & 0.37 & 0.41 \\
\tablelang{Action age at execution (robot)}{执行时的动作年龄（真机）} & \resultplaceholder & \resultplaceholder \\
\tablelang{Queries over 50\,ms (robot, \%)}{超过 50\,ms 的查询（真机，\%）} & \multicolumn{2}{c}{\resultplaceholder} \\
\bottomrule
\end{tabular}

\end{table}
\end{minipage}

\subsection{Results on DOMINO}
\label{app:external}
The fixed W5 split covers 24 clean level-1 Franka left-arm tasks from DOMINO
\cite{domino}. Each task contributes 50 development and 100 test scenes, for
1{,}200 development and 2{,}400 test scenes. All four methods use the same test
scenes, held out from training, validation, and rollout selection. The native
evaluator pauses physics during inference. The manipulation score (MS) scales
route completion by penalties for leaving the workspace or view and for
collisions. \ours reads the ground-truth pose of the instructed target from the
simulator and masks slots without a valid pose (Appendix~\ref{app:tokens}). The
RGB-based baselines do not receive simulator target poses.

\paragraph{Data and adaptation.}
Our C12 model uses the catalog-v4 release of 1{,}200 training and 240 validation
demonstrations. The baselines use the earlier DOMINO Franka corpus with the same
nominal split sizes. Policies execute absolute joint-position targets with DOMINO's
native controller in its eight-dimensional action space. Normalization is fitted
on training data. The official DynamicWAM sampler skips demonstrations too short
for its windows, leaving 1{,}054 training and 210 validation demonstrations and
none for two tasks. The test set retains all 24 tasks. DynamicWAM starts from its
published checkpoint and trains for 40{,}000 updates at global batch 128.
$\pi_{0.5}$ starts from its official base and trains for 30{,}000 updates at global
batch 256. Evaluation uses its final EMA. PUMA starts from its published weights,
keeps its teacher frozen, and trains at global batch 64. Validation action MSE
selects update 6{,}000. \ours trains from scratch with global batch 512 and stops
after 1{,}575 updates. The lowest validation loss selects update 1{,}365.
Development rollouts select three actions per query from its eight-action chunk.
At test time, \ours, $\pi_{0.5}$, DynamicWAM, and PUMA execute 3, 10, 16, and 16
actions per query, respectively.

\subsection{Ablation Controls}
\label{app:ablations}
\label{app:ablation_defs}
The six controls of Sec.~\mainref{sec:principles} (Table~\mainref{tab:analysis}) share the
data, training recipe, and executor. Each is evaluated at the final checkpoint of its
run. Removing a loss also removes its head. The control without target history removes
slots 0--6 and keeps slot 7, the current pose (Appendix~\ref{app:tokens}). The uniform
fine flow control removes $\mathcal L_P$ and the pose head. It supervises flow at
$4\times4$ points in each patch, forming a $56\times56$ grid with labels from the same
ground-truth rigid motion.
